\scanpage{30}
\mspage{29}
% Reused typed sheet, read after rotation through 180 degrees.
\begin{center}-2.17-\end{center}
\[
\begin{aligned}
\lambda(\alpha)\cdot\gamma\cdot\lambda(\alpha^{-1})
&=\chi(\lambda(\alpha^{-1}))\cdot\gamma\\
&=\alpha^{-\bar\chi(\bar\lambda)}\cdot\gamma
\end{aligned}
\]
for $\gamma\in(\mathbb G_a)_k$. Therefore $\lambda(\alpha)\cdot\gamma\cdot\lambda(\alpha^{-1})$ has a finite specialization when $\alpha\longrightarrow0$, i.e. $\mathbb G_a\subset P(\lambda)$. Therefore $B\subset P(\lambda)$. \uline{QED}

\uline{Proposition 2.7}: Let a reductive group $G$ act via $\sigma$ on a scheme $X$, proper over $k$. Then for all $x\in X_k$, $L\in\mathrm{Pic}^G(X)$, and 1-PS's $\lambda$ of $G$,
\[
\mu^L(x,\lambda)=\mu^L(x,\gamma^{-1}\cdot\lambda\cdot\gamma)
\]
if $\gamma\in P(\lambda)_k$.

\uline{Proof}: Let $y$ be the specialization of $\sigma(\lambda(\alpha),x)$ as $\alpha\longrightarrow0$. Suppose we write $y=\lim_{\alpha\to0}\sigma(\lambda(\alpha),x)$. In this notation, we calculate:
\[
\begin{aligned}
&\lim_{\alpha\to0}\sigma(\gamma^{-1}\cdot\lambda(\alpha)\cdot\gamma,x)\\
&\quad=\lim_{\alpha\to0}\sigma\bigl(\gamma^{-1}\cdot\lambda(\alpha)\cdot\gamma\cdot\lambda(\alpha^{-1}),\sigma(\lambda(\alpha),x)\bigr)\\
&\quad=\sigma\bigl(\gamma^{-1}\cdot\lim_{\alpha\to0}\lambda(\alpha)\cdot\gamma\cdot\lambda(\alpha^{-1}),y\bigr).
\end{aligned}
\]
Let $\lim_{\alpha\to0}\lambda(\alpha)\cdot\gamma\cdot\lambda(\alpha^{-1})=\gamma'$. According to Proposition 2.6, $\gamma'$ centralizes $\lambda$. Now $\mu(x,\gamma^{-1}\cdot\lambda\cdot\gamma)$ is calculated by looking at $\lim_{\alpha\to0}\sigma(\gamma^{-1}\cdot\lambda(\alpha)\cdot\gamma,x)$. Therefore:
\[
\mu(x,\gamma^{-1}\cdot\lambda\cdot\gamma)
=\mu(\sigma(\gamma^{-1}\cdot\gamma',y),\gamma^{-1}\cdot\lambda\cdot\gamma)
\]
By property (i) of $\mu$ (\S1 above), we have:
